% HAND-WRITTEN fragment -- NOT generated by maestro2tex. Input from
% AnalogChapter.tex after note_eis_accuracy.tex.
%
% Sources (~/chips/castalia/simruns/eisdemo_20260825):
%   coverage_map.txt (README S5) -- the table below, verbatim:
%     ncode Rf_meas delta%  |Z|min  |Z|max  span  f_Cf
%     0     19630.4 +9.06   236.01  150.87k 639.3  nan (>100 kHz sweep; 1/(2 pi
%                                                   19.63k 22p) = 368 kHz)
%     20   138722.1 +0.52   1667.8  1.0662M 639.3  51.98k
%     63   394838.7 -0.29   4747    3.0346M 639.3  18.21k
%     102  628430.5 -0.25   7555.4  4.8298M 639.3  11.46k
%     191 1159552.3 -0.38   13941   8.9118M 639.3   6.238k
%     319 1924273.3 -0.40   23135   14.789M 639.3   3.746k
%   Union 236.01 ohm - 14.789 MOhm = 4.80 decades.  |Z|min is where the square's
%   PEAK current fills the +-0.80 V compliance; |Z|max where the FUNDAMENTAL of
%   the output falls to one converter LSB (1.5934 mV).  Excitation is the
%   MEASURED delivered square, 19.2362 mV pp = 9.618 mV peak = 12.246 mV
%   fundamental.  f_Cf is the -3 dB frequency of the measured |Zt|, i.e. the real
%   22 pF feedback pole: 52.0 kHz measured vs 52.2 predicted at N = 20, 3.746 vs
%   3.76 kHz at N = 319.
%   ratio/band_edges.txt (README S4) -- each gain code against a pure resistor of
%   exactly Rf/ratio, so the noise gain is frozen across the sweep:
%     err(dc) at N = 0: 0.1204 / 0.3304 / 0.6603 / 1.3217 % at ratio 7.3/20/40/80
%     ratio 40: f(+-1 %)   22.67k (N0) 14.13k (N20) 9.849k (N63) 9.102k (N102)
%                          13.87k (N191) >100k (N319)
%               f(0.573 deg) 5527 (N0) 5760 (N20) 6930 (N63) >100k (N102+)
%   Predicted edges, afe_rev2_eis_lite.html S5 at ratio 40: 5.56 kHz for both.
%   Model grading (README S3.4(1)): reproduces the measured error to <= 7 % at
%   noise gain >= 13; 42 % optimistic at noise gain 2.8 (0.0294 % measured vs
%   0.0207 % model).
% Library cell names are kept out of the rendered prose.

\subsubsection{Coverage and Bandwidth}\label{sss:eis-coverage}

At the measured excitation the channel covers \SI{236}{\ohm} to
\SI{14.79}{\mega\ohm}, \num{4.80} decades (Table~\ref{tab:eis-coverage}). Each
gain code spans \(639\times\), \num{2.81} decades, between the impedance at
which the excitation's peak current fills the \(\pm\SI{0.80}{\volt}\) compliance
and the impedance at which the fundamental of the output falls to one converter
code. Adjacent codes overlap by more than a decade, so \textbf{there are no
impedance seams}: the 320-tap ladder tiles that span as a continuum, and the
only seams in the sweep are the ladder's own carry codes
(Section~\ref{ss:tiarprog}), which move a window edge by far less than the
overlap. The frequency and impedance limits do not collide either --- a
capacitive interface's impedance falls as frequency rises, so a sweep walks down
the table as it walks up in frequency, and the high-gain codes whose feedback
pole sits between \num{3.7} and \SI{11}{\kilo\hertz} are only selected at the
low-frequency end. That pole is measured here and confirms the
\SI{22}{\pico\farad} compensation capacitor to better than \SI{1}{\percent}:
\SI{52.0}{\kilo\hertz} against \SI{52.2}{\kilo\hertz} predicted at gain code 20,
\SI{3.75}{\kilo\hertz} against \SI{3.76}{\kilo\hertz} at code 319.

\paragraph{The band edge is a phase specification, not a magnitude one.}
Presenting each gain code with a pure resistor of exactly \(R_f/40\) freezes the
noise gain across the sweep and measures the edge directly. The dc error then
scales exactly with that ratio --- \num{0.120}, \num{0.330}, \num{0.660},
\SI{1.322}{\percent} at \num{7.3}, \num{20}, \num{40}, \num{80} --- so the form
of the design equation holds. Its bandwidth prediction does not: the
\(\pm\SI{1}{\percent}\) magnitude edge measures \SIrange{9.1}{22.7}{\kilo\hertz}
across the gain codes against \SI{5.56}{\kilo\hertz} predicted, \num{1.6} to
\(4\times\) pessimistic, while the \(\pm\SI{0.573}{\degree}\) phase edge measures
\SIrange{5.53}{6.93}{\kilo\hertz} --- the predicted value. \textbf{The usable
band is therefore \SI{0.1}{\hertz} to \SI{5.5}{\kilo\hertz} set by phase}, with
magnitude in specification to at least \SI{9.1}{\kilo\hertz}. The same model
under-predicts the error itself by \SI{42}{\percent} where the noise gain is
near unity (\num{0.029} measured against \(0.021\,\%\)), which at that size
changes nothing.

\begin{table}[htbp]
  \centering
  \caption[{Impedance window and feedback pole against gain code.}]{Impedance
  window and measured feedback pole against gain code, at the delivered
  \SI{19.24}{\milli\volt} excitation. \(R_f\) is the transimpedance measured in
  the closed loop at \SI{0.1}{\hertz}, compared with the ladder's nominal law.
  The window runs from the impedance at which the excitation's peak current fills
  the \(\pm\SI{0.80}{\volt}\) output compliance to the impedance at which the
  fundamental of the output falls to one converter code; the peak sets the lower
  edge because the peak is what rails, and the fundamental sets the upper edge
  because the fundamental is what the demodulator reads. Adjacent windows overlap
  by more than a decade. The last column is the \SI{-3}{\decibel} frequency of
  the measured transimpedance, i.e. the real \SI{22}{\pico\farad} feedback pole.}
  \label{tab:eis-coverage}
  \begin{tabular}{rrrrrr}
    \toprule
    $N$ & $R_f$ measured & vs law & $|Z|_{\min}$ & $|Z|_{\max}$ & Feedback pole \\
    \midrule
    0   & \SI{19.63}{\kilo\ohm} & \(+9.06\,\%\) & \SI{236}{\ohm}       & \SI{150.9}{\kilo\ohm} & \(>\SI{100}{\kilo\hertz}\)$^{a}$ \\
    20  & \SI{138.7}{\kilo\ohm} & \(+0.52\,\%\) & \SI{1.67}{\kilo\ohm} & \SI{1.066}{\mega\ohm} & \SI{52.0}{\kilo\hertz} \\
    63  & \SI{394.8}{\kilo\ohm} & \(-0.29\,\%\) & \SI{4.75}{\kilo\ohm} & \SI{3.035}{\mega\ohm} & \SI{18.2}{\kilo\hertz} \\
    102 & \SI{628.4}{\kilo\ohm} & \(-0.25\,\%\) & \SI{7.56}{\kilo\ohm} & \SI{4.830}{\mega\ohm} & \SI{11.5}{\kilo\hertz} \\
    191 & \SI{1.160}{\mega\ohm} & \(-0.38\,\%\) & \SI{13.9}{\kilo\ohm} & \SI{8.912}{\mega\ohm} & \SI{6.24}{\kilo\hertz} \\
    319 & \SI{1.924}{\mega\ohm} & \(-0.40\,\%\) & \SI{23.1}{\kilo\ohm} & \SI{14.79}{\mega\ohm} & \SI{3.75}{\kilo\hertz} \\
    \bottomrule
    \multicolumn{6}{l}{\footnotesize $^{a}$\,Outside the \SI{100}{\kilo\hertz}
      sweep; the element values put it at \SI{368}{\kilo\hertz}.} \\
  \end{tabular}
\end{table}
